\section{Python Code to RDPro Code}
In this section, we use NDVI computation as a concrete example to illustrate  the translation from single machine Python to distributed RDPro Scala code. NDVI (Normalized Difference Vegetation Index) is a widely used index in agriculture and remote sensing, defined as:
\begin{equation}
    NDVI = \frac{NIR - Red}{NIR + Red}
\label{eq:ndvi}
\end{equation}

\noindent where $NIR$ and $Red$ denote the pixel values of the near-infrared and red spectral bands, respectively. 

Figure~\ref{fig:ndvi-comparison} shows the two implementations side by side. The input raster files are Landsat8 satellite images where band values are 
stored as 16-bit unsigned integers. The input Python code loads both bands as 32-bit floating point arrays into memory on a single machine using GDAL, constructs a boolean NoData mask, and applies the NDVI formula. The RDPro implementation generated by the LLM represents each band as a distributed \texttt{RasterRDD[Float]}. Rather than combining bands through direct 
array arithmetic, the two RDDs must first be spatially aligned via the 
\texttt{overlay()} operation, producing a \texttt{RasterRDD[Array[Float]]}. 
The NDVI formula is then applied per-pixel using a \texttt{mapPixels\{\}} 
lambda. While the generated RDPro code is logically correct, it fails to 
handle data types consistently because the Python code explicitly casts both bands to \texttt{float32} upon reading, whereas RDPro does not apply an equivalent cast, causing a type mismatch error at runtime. This subtle but critical difference illustrates a common failure mode in LLM-generated translations, where the model preserves the high-level logic but overlooks type system constraints specific to the target framework.
To address these failure modes, we introduce three modifications to the 
RDPro codebase that improve the reliability of LLM-generated code.

% \begin{figure*}[t]
% \begin{minipage}{0.48\linewidth}
% \lstset{style=codeStyle, language=Python}
% \begin{lstlisting}
% from osgeo import gdal
% import numpy as np

% # Load Red (B4) and NIR (B5) bands
% ds_red = gdal.Open(B4_PATH, gdal.GA_ReadOnly)
% ds_nir = gdal.Open(B5_PATH, gdal.GA_ReadOnly)

% red = ds_red.GetRasterBand(1).ReadAsArray()
%         .astype(np.float32)
% nir = ds_nir.GetRasterBand(1).ReadAsArray()
%         .astype(np.float32)

% # NoData mask
% mask = (red == red_nodata) | (nir == nir_nodata)

% # NDVI formula
% denom = nir + red
% ndvi = np.where((denom == 0) | mask,
%                 -9999.0,
%                 (nir - red) / denom)

% # Write output GeoTIFF
% driver = gdal.GetDriverByName("GTiff")
% out_ds = driver.Create(OUT_NDVI,
%              ds_red.RasterXSize,
%              ds_red.RasterYSize, 1,
%              gdal.GDT_Float32,
%              options=["TILED=YES", "COMPRESS=LZW"])
% out_ds.GetRasterBand(1).WriteArray(ndvi)
% \end{lstlisting}
% \end{minipage}
% \hfill
% \begin{minipage}{0.48\linewidth}
% \lstset{style=codeStyle, language=Scala}
% \begin{lstlisting}
% import edu.ucr.cs.bdlab.beast._
% import edu.ucr.cs.bdlab.raptor._

% // Load Red (B4) and NIR (B5) as distributed RDDs
% val redBand: RasterRDD[Float] =
%     sc.geoTiff(b4Path)
% val nirBand: RasterRDD[Float] =
%     sc.geoTiff(b5Path)

% // Spatially align bands into a single RDD
% val stacked: RasterRDD[Array[Float]] =
%     redBand.overlay(nirBand)

% // NDVI formula applied per pixel
% val ndvi: RasterRDD[Float] =
%     stacked.mapPixels { px =>
%       val (red, nir) = (px(0), px(1))
%       val denom = nir + red
%       if (denom == 0f || isNoData(red, nir))
%         -9999.0f
%       else
%         (nir - red) / denom}

% // Write output GeoTIFF
% ndvi.saveAsGeoTiff(outPath,
%          Seq(
%           GeoTiffWriter.Compression -> TiffConstants.COMPRESSION_LZW,
%           GeoTiffWriter.WriteMode -> "compatibility"
%         ))
% \end{lstlisting}
% \end{minipage}
% \caption{NDVI computation from a Landsat8. \textbf{Left:} Python 
% script using GDAL. \textbf{Right:} RDPro implementation by LLM in Scala.}
% \label{fig:ndvi-comparison}
% \end{figure*}